\documentclass[10.5pt,a4paper]{article}
\usepackage[margin=2cm]{geometry}
\usepackage[T1]{fontenc}
\usepackage{lmodern,amsmath,amssymb,longtable,booktabs,array,enumitem}
\usepackage[hidelinks]{hyperref}
\renewcommand{\arraystretch}{1.2}\setlength{\tabcolsep}{4pt}
\newcolumntype{L}[1]{>{\raggedright\arraybackslash}p{#1}}
\newcommand{\cls}[1]{\textsf{#1}}
\title{CRASH: Collision Risk Analysis through Semantic Hierarchies\\[4pt]\large AI-3001 Iteration 1 --- Part B: Ontology / Knowledge Model}
\author{Team members: Ritaj, Hamna, Tahreem}\date{November 2026}
\begin{document}\maketitle

\section{Overview}
The model is an OWL 2 ontology in one Turtle file (\texttt{crash\_schema.ttl}), built in Protégé. Namespace: \texttt{https://crash.example.org/ontology\#}. It has \textbf{17 classes, 5 object properties and 17 data properties}, and five reasoning constructs (the guideline asks for 10, 8 and 3). Data: STATS19 2025, all collisions of North Wales and Gwent police (1,115 collisions).

\section{T-Box and A-Box}
\begin{itemize}[nosep]
\item \textbf{T-Box (schema)}: the classes, class hierarchy, properties, domain/range, functional and disjointness constraints, and the defined classes. It is the file \texttt{crash\_schema.ttl}.
\item \textbf{A-Box (instances)}: the facts from the CSVs. One individual per collision, vehicle and casualty; one per LSOA, local authority and police force. Example IRIs: \texttt{Collision\_202560C131039}, \texttt{LSOA\_W01000040}, \texttt{PoliceForce\_60}. The A-Box is built from the CSVs in the next step.
\end{itemize}

\section{Class hierarchy}
\begin{verbatim}
Collision                     Casualty                      Area
  FatalCollision                VulnerableRoadUser            LSOA
  SeriousCollision                Pedestrian                    HighRiskArea
Vehicle                           Cyclist                     LocalAuthority
                                  Motorcyclist                PoliceForce
                                ChildCasualty

Defined classes (inferred, not stored):
  KSICollision = FatalCollision or SeriousCollision
  VRUCollision = Collision with involvesCasualty some VulnerableRoadUser
\end{verbatim}
KSICollision and VRUCollision are \emph{defined} classes: no collision is stored as one, the reasoner infers membership.

\section{Properties}
\small
\begin{longtable}{L{3.4cm} L{2.6cm} L{2.6cm} L{7.2cm}}
\toprule \textbf{Object property} & \textbf{Domain} & \textbf{Range} & \textbf{Characteristic}\\\midrule
hasVehicle & Collision & Vehicle &\\
hasCasualty & Vehicle & Casualty &\\
involvesCasualty & Collision & Casualty & Chain: hasVehicle $\circ$ hasCasualty\\
occurredIn & Collision & LSOA & Functional\\
partOf & Area & Area & Transitive (LSOA $\to$ LocalAuthority $\to$ PoliceForce)\\\midrule
\textbf{Data property} & \textbf{Domain} & \textbf{Range} & \textbf{Meaning}\\\midrule
collisionIndex, collisionDate, speedLimit & Collision & string, date, integer & Key, date, mph\\
roadClass, roadType, junctionType & Collision & string & Text labels from the DfT data guide\\
lightCondition, weatherCondition, surfaceCondition & Collision & string & e.g. ``Daylight'', ``Wet or damp''\\
vehicleType, vehicleManoeuvre, driverAge & Vehicle & string, string, integer & driverAge omitted if missing\\
casualtySeverity, casualtyClass, casualtyAgeBand & Casualty & string, string, integer & Fatal/Serious/Slight; band 1--3 = child\\
areaCode, areaName & Area & string & LSOA / ONS code, name\\
\bottomrule\end{longtable}
\normalsize
All data properties are functional. Codes $-1$ and 99 become ``Unknown''; flag and adjusted-severity columns are dropped.

\section{Constraints}
Domain and range on every property; functional properties (one LSOA, one road type per collision); disjoint classes: Collision, Vehicle, Casualty and Area; LSOA, LocalAuthority and PoliceForce; FatalCollision and SeriousCollision.

\section{Reasoning constructs and rules}
\small
\begin{longtable}{L{0.8cm} L{6.2cm} L{4.2cm} L{4.8cm}}
\toprule \textbf{ID} & \textbf{Construct} & \textbf{Where} & \textbf{Inferred fact}\\\midrule
R1 & KSICollision $\equiv$ FatalCollision $\sqcup$ SeriousCollision & Defined class & Serious collision is a KSICollision\\
R3 & VRUCollision $\equiv$ $\exists$involvesCasualty.VulnerableRoadUser & Defined class & Collision with a cyclist casualty\\
R4 & Pedestrian, Cyclist, Motorcyclist $\sqsubseteq$ VulnerableRoadUser $\sqsubseteq$ Casualty & Subclass hierarchy & A motorcyclist is a Casualty\\
R5 & partOf transitive & Object property & LSOA partOf PoliceForce\\
R6 & involvesCasualty $\equiv$ hasVehicle $\circ$ hasCasualty & Property chain & Collision involvesCasualty $c$\\
R7 & LSOA with $\geq 3$ KSICollisions $\Rightarrow$ HighRiskArea & SPARQL INSERT (GraphDB) & 30 of 510 LSOAs\\
R8 & Age band 1--3 $\Rightarrow$ ChildCasualty & SPARQL INSERT (GraphDB) & Casualty aged 12\\
\bottomrule\end{longtable}
\normalsize
R1--R6 run in OWL2-RL; GraphDB has no SWRL, so the counting rule R7 and the number-range rule R8 are SPARQL INSERT queries (text in the file header).

\section{Competency questions covered}
\small
\begin{tabular}{L{2.3cm} L{13.7cm}}
\toprule CQ1, CQ13, CQ14 & Collision, Vehicle, Casualty, VulnerableRoadUser, VRUCollision, involvesCasualty (R3, R4, R6)\\
CQ2, CQ5, CQ8, CQ11 & KSICollision with roadClass, surfaceCondition, lightCondition, roadType, junctionType, occurredIn, partOf (R1, R5)\\
CQ3, CQ13 & LSOA, LocalAuthority, PoliceForce, HighRiskArea (R5, R7)\\
CQ4, CQ7, CQ10 & Pedestrian, Cyclist, Motorcyclist, casualtySeverity, vehicleManoeuvre, speedLimit (R4, R6)\\
CQ6 & KSICollision, VRUCollision, partOf (R1, R3, R5, R6)\\
CQ9 & ChildCasualty (R8)\\
CQ12 & collisionIndex and all properties\\
\bottomrule\end{tabular}
\normalsize

\section{Check in Protégé}
Open the file, start HermiT. A test collision typed SeriousCollision with a vehicle that has a motorcyclist casualty was inferred to be a KSICollision and a VRUCollision, with involvesCasualty to the casualty (screenshots kept for the review). The six test individuals are then deleted (next section) so the file holds only the schema.

\section{Removing the test instances}
In Protégé: Individuals tab $\rightarrow$ select \texttt{c1}, \texttt{v1}, \texttt{k1}, \texttt{lsoa1}, \texttt{la1}, \texttt{pf1} (Ctrl-click) $\rightarrow$ delete button (purple diamond with a cross) $\rightarrow$ confirm. Stop the reasoner, then File $\rightarrow$ Save as Turtle, to keep the A-Box empty and the T-Box unchanged.
\end{document}
